% Figure: a car parked on a slope, FBD with axle normals, weight, and
% the friction holding it. Used in the worked example and exercise 1.6.
\begin{figure}[htbp]
\centering
\begin{tikzpicture}[
  wt/.style={draw=engblue, -{Stealth}, very thick},
  nrm/.style={draw=engred, -{Stealth}, very thick},
  fric/.style={draw=engamber, -{Stealth}, very thick},
  scale=1.0,
]
% incline at ~12 deg for visibility
\begin{scope}[rotate=12]
  % road surface
  \draw[draw=engblack, very thick] (-0.6,0) -- (7.4,0);
  \foreach \x in {-0.3,0.2,...,7.2} {\draw[draw=enggray, thin] (\x,0) -- (\x-0.2,-0.25);}
  % car body (simplified)
  \draw[draw=engblack, fill=engblue!10, thick]
    (1.3,0.55) -- (2.0,1.25) -- (4.6,1.25) -- (5.4,0.55) -- cycle;
  \draw[draw=engblack, fill=white, thick] (2.2,0.55) circle (0.45);
  \draw[draw=engblack, fill=white, thick] (4.5,0.55) circle (0.45);
  \fill[engblack] (2.2,0.55) circle (0.05);
  \fill[engblack] (4.5,0.55) circle (0.05);
  % centre of mass
  \coordinate (G) at (3.6,1.0);
  \fill[engblack] (G) circle (0.06);
  \node[font=\scriptsize, anchor=south] at (G) {$G$};
  % normals at each axle (perpendicular to road = vertical in this frame)
  \draw[nrm] (2.2,0.1) -- ++(0,1.6) node[anchor=south, font=\scriptsize, engred] {$N_r$};
  \draw[nrm] (4.5,0.1) -- ++(0,1.6) node[anchor=south, font=\scriptsize, engred] {$N_f$};
  % friction up the slope
  \draw[fric] (2.2,0.1) -- ++(-1.6,0) node[anchor=east, font=\scriptsize, engamber] {$F_f$};
  % weight straight down (true vertical = -90-12 deg in the rotated frame)
  \draw[wt] (3.6,1.0) -- ++(-102:2.0) node[anchor=north, font=\scriptsize, engblue] {$mg$};
\end{scope}
% slope angle marker at origin
\draw[draw=enggray, thin] (1.2,0) arc (0:12:1.2);
\node[font=\scriptsize] at (1.45,0.12) {$\theta$};
\end{tikzpicture}
\caption[Free-body diagram of a car parked on a slope]{A parked car on
a slope of angle $\theta$. The chosen body is the whole car. Each axle
contact contributes a normal force perpendicular to the road, $N_f$ at
the front and $N_r$ at the rear; the braked or chocked wheels supply a
friction force $F_f$ up the slope; the weight $mg$ acts straight down at
the centre of mass $G$. Resolving along and perpendicular to the road
gives $F_f = mg\sin\theta$ and $N_f + N_r = mg\cos\theta$, while moments
about one axle split the normal load between front and rear according to
where $G$ sits in the wheelbase. The steeper the slope, the more of the
weight the friction must hold and the more the normal load shifts toward
the downhill axle.}
\label{fig:vol03ch01:parked-car}
\end{figure}
